\documentclass{beamer}
\usetheme[block=fill]{metropolis}
\usepackage{graphicx}
\graphicspath{{../../latex/images/}}
\usepackage{boolexpr}
\usepackage{kpfonts}
\usepackage{caption}
\usepackage{wrapfig}
\usepackage[labelformat=empty]{caption}
\usepackage[official]{eurosym}

\usepackage[charsperline=60]{jlcode}
\addlitjlmacros{@constraint}{@constraint}{11}
\addlitjlmacros{@variable}{@variable}{9}
\addlitjlmacros{@objective}{@objective}{10}

\usepackage{minted}

\usepackage{siunitx}

\usepackage[math]{iwona}
\usepackage{scalerel}
\def\lVert{\mid\!\mid}
\def\rVert{\mid\!\mid}

\usepackage[normalem]{ulem}
\usepackage{mathtools}

\usepackage{../../latex/custom}
\usepackage{amsfonts}

\newcommand{\expe}[2]{\la #1, #2 \ra}

\usepackage{framed}

\newcommand\Wider[2][3em]{%
\makebox[\linewidth][c]{%
  \begin{minipage}{\dimexpr\textwidth+#1\relax}
  \raggedright#2
  \end{minipage}%
  }%
}

\newcommand{\aeur}{\alpha_\text{\euro}}
\newcommand{\adol}{\alpha_\$}
\newcommand{\apou}{\alpha_\text{\pounds}}

\title{Multivariate polynomials in Julia}
\date{July, 2022}
\author{Beno\^it Legat}
\institute{
  Joint work with Robin Deits, Chris Elrod and Yingbo Ma\\
  JuliaCon-dev 2022 \hfill \includegraphics[scale=0.2]{juliacon.png}\\[5em]
}

\makeatletter
\let\@@magyar@captionfix\relax
\makeatother
\begin{document}
  \maketitle

  \begin{frame}
    \frametitle{Where are multivariate polynomials used in Julia ?}
    Used in variety of applications:
    \begin{itemize}
      \item HomotopyContinuation: Solve polynomial systems with homotopy continuation
      \item SumOfSquares: Sum-of-Squares programming
      \item Symbolics: Manipulate and simplify symbolic expressions
      \item Groebner: Gröbner bases computations based on Faugère's F4 algorithm
      \item SemialgebraicSets: Reduction modulo ideal with Gröbner bases + solve polynomial systems with Schur decomposition
      \item Nemo/AbstractAlgebra: Computer algebra system
      \item CondensedMatterSOS: Condensed Matter with Sum-of-Squares
      \item ...
    \end{itemize}
  \end{frame}

  \begin{frame}[fragile]
    \frametitle{How are multivariate polynomials represented?}
    Sorted vector of terms:
\begin{jllisting}
struct Term{C,M}
    coefficient::C
    monomial::M
end
struct Polynomial{C,M}
    terms::Vector{Term{T,M}}
end
\end{jllisting}

    \uncover<2->{
      Sorted how ? \alert{Any} ordering ?
    }

    \uncover<3->{
      If $m_1 > m_2$ and $m_3 > m_4$, we need $m_1m_3 > m_2m_4$.\\
      {\scriptsize See 2.\S2.Definition 1 of Cox Little O'Shea ``\emph{Ideal, Varieties, and Algorithms}''}.
    }
  \end{frame}

  \begin{frame}[fragile]
    \frametitle{How are monomials represented?}
    \begin{block}{DynamicPolynomials}
      \vspace{-0.5cm}
\begin{jllisting}
struct Variable
    id::Int
    name::String
end
struct Monomial
    variables::Vector{Variable}
    exponents::Vector{Int}
end
\end{jllisting}
      \vspace{-0.5cm}
    \end{block}

    \begin{block}{TypedPolynomials (Robin Deits)}
      \vspace{-0.5cm}
\begin{jllisting}
struct Variable{Name} end
struct Monomial{V<:Tuple{Vararg{Variable,N}},N}
    exponents::NTuple{N,Int}
end
\end{jllisting}
    \end{block}
      \vspace{-1cm}
  \end{frame}

  \begin{frame}[fragile]
    \frametitle{How are monomials represented? (part 2)}
    \begin{block}{SIMDPolynomials (Chris Elrod and Yingbo Ma)}
      Bit packed with maximum $L$ variables and $E$ bits of exponents (so $64K > LE$).
\begin{jllisting}
struct Monomial{L,E,K}
    bits::NTuple{K,UInt64}
end
\end{jllisting}
    \end{block}
    \begin{block}{CondensedMatterSOS (Jacopo Surace)}
\begin{jllisting}
struct Monomial
    variables::SortedDict{Int,SpinVariable}
end
\end{jllisting}
    \end{block}

  \end{frame}

  \begin{frame}[fragile]
    \frametitle{MultivariatePolynomials API}
    \begin{itemize}
      \item \jlinl{AbstractVariable}, \jlinl{AbstractMonomial}, \jlinl{AbstractTerm}, \jlinl{AbstractPolynomials}
      \item Implementation-independent API: \jlinl{degree}, \jlinl{coefficient}, \jlinl{monomial}, \jlinl{isless}, \jlinl{+}, ...
      \item Implements MutableArithmetics to make \alert{generic} code \alert{fast}
      \item Pretty printing
      \item Differentiation
      \item Substitution (partial) and evaluation (i.e., full substitution)
      \item Euclidean division and \alert{gcd}/\alert{gcdx}
    \end{itemize}
    Get all that for free for any new fancy monomial representation!
  \end{frame}

\end{document}
